\documentclass[aps,onecolumn,nofootinbib,pra]{article}

\usepackage{../../spec_files/arxiv}
\usepackage{amsmath,amsfonts,amssymb,amsthm,enumerate,times}
\usepackage{mathtools}
\usepackage[usenames,dvipsnames]{color}
\usepackage{hyperref}
\hypersetup{
    breaklinks,
    colorlinks,
    linkcolor=gray,
    citecolor=gray,
    urlcolor=gray,
    pdftitle={},
    pdfauthor={Alex Goessmann}
}

\usepackage{tikz}
\usepackage{graphicx}
\usepackage{float}
\usepackage{comment}
\usepackage{csquotes}
\usepackage{braket}
\usepackage{listings}
\usepackage{verbatim}
\usepackage{etoolbox}
\usepackage[utf8]{inputenc}
\usepackage[english]{babel}
\usepackage[T1]{fontenc}
\usepackage{titlesec}
\usepackage{fancyhdr}
\usepackage{bbm}
\usepackage{bm}
\usepackage{algpseudocode}
\usepackage{algorithm}
\usepackage{lipsum}

\newtheorem{remark}{Remark}
\newtheorem{theorem}{Theorem}
\newtheorem{lemma}{Lemma}
\newtheorem{corollary}{Corollary}
\newtheorem{definition}{Definition}
\newtheorem{example}{Example}

\pretolerance=500
\tolerance=100
\emergencystretch=10pt

% Bibliography
\DeclareUnicodeCharacter{FB01}{fi}
\usepackage[round]{natbib}
\usepackage{wasysym}

\newcommand{\red}[1]{\textcolor{red}{#1}}

\usepackage{minted}

\usepackage{import}
\subimport{../../macros/}{all-macros}

\begin{document}
    \title{Bayesian inversion for PDE models in \tnreason{}}

    \maketitle
    \date{\today}

    \begin{abstract}
        We represent all solutions of a PDE to different parameters as tensor networks.
        This enables the integration of logical constraints and formulation of inference tasks, such as sampling from Bayesian posteriors, based on contractions.
    \end{abstract}


    \section{PDE solutions as tensors}

    Given a discretization of a function, we investigate the indicator tensor, whether a combination of values is a solution to the PDE.

    \subsection{Tensor representation of solutions}

    First of all we discretize continuous functions towards a tensor representation:
    \begin{itemize}
        \item Selection of $\catorder$ domain elements
        \item Selection of $\catdim$ image elements together with an approximator mapping each image element to one of the selected
        \item Variables $\catvariableof{\catenumerator}$ representing the function evaluation at the $\catenumeratorin$ domain element
    \end{itemize}

    In this discretization scheme any functions is represented by an index $\shortcatindices$.
    A set of function is represented by its basis encoding
    \begin{align*}
        \hypercoreofat{\mathcal{F}}{\indexedshortcatvariables}
        = \begin{cases}
              1 & \ifspace \shortcatindices \in \mathcal{F} \\
              0 & \elsetext \, .
        \end{cases}
    \end{align*}

    We are interested in the discretized solution sets $\mathcal{F}$ of PDEs and in decompositions of $\hypercoreofat{\mathcal{F}}{\shortcatvariables}$.
    Here, we include parameters of the PDE as further variables $\thirdcatvariable$.

    \begin{example}[One-dimensional Darcy flows]
        In the example of one-dimensional Darcy flows, the variables $\catvariableof{\catenumerator}$ describe the pressure at a position $\catenumerator$ and further variables $\seccatvariableof{\catenumerator}$ the medium resistance at that position.
    \end{example}

    \subsection{Euler method}

    With Eulers method a solution to a PDE is approximated step by step.
    Given the value $\catindexof{\catenumerator}$ the $\catindexof{\catenumerator+1}$ is computed given the parameters $\thirdcatindex$.
    We define $\formulaof{\catenumerator}$ as the function mapping $\catindexof{\catenumerator}$ and $\thirdcatindex$ to the corresponding $\catindexof{\catenumerator+1}$.

    The set of Euler approximations $\mathcal{F}$ of the solutions to a PDE is then decomposed as
    \begin{align*}
        \hypercoreofat{\mathcal{F}}{\shortcatvariables,\thirdcatvariable}
        = \contractionof{
            \{\bencodingofat{\formulaof{\catenumerator}}{\catvariableof{\catenumerator+1},\catvariableof{\catenumerator},\thirdcatvariable}
            \wcols \catenumerator\in[\catorder-1]\}
        }{\shortcatvariables,\thirdcatvariable} \, .
    \end{align*}

    \begin{example}[Toy example]\label{exa:EulerToy}
        Consider the parameter free toy example $u'=1$ on the interval $[0,1]$ discretized into $\catorder$ equal size cells.
        We choose $\catdim\geq \catorder$ and enumerate $[\catdim]$ another discretization of an interval with equal cell size.
        The discretized solution is then captured by the functions
        \begin{align*}
            \formulaof{\catenumerator}(\catindexof{\catenumerator})
            = \catindexof{\catenumerator}+1
        \end{align*}
        and the solutions tensor is the decomposition:

        \stantikzpic{
            \drawhblock{0}{0}{2}{$\bencodingof{\formulaof{0}}$}
            \drawtdwire{-1}{-1}{$\catvariableof{0}$}
            \drawddwire{1}{-3}{$\catvariableof{1}$}

            \drawvariabledot{2.5}{-5}
            \draw[-<-] (2.5,-5) to [bend right= -10] (1,-3);
            \draw[->-] (2.5,-5) to [bend right= 10] (4,-3);
            \drawddwire{2.5}{-7}{$\catvariableof{1}$}

            \drawhblock{5}{0}{2}{$\bencodingof{\formulaof{1}}$}
            \drawtdwire{4}{-1}{$\catvariableof{1}$}
            \drawddwire{6}{-3}{$\catvariableof{2}$}


            \begin{scope}[shift={(5,0)}]
                \drawvariabledot{2.5}{-5}
                \draw[-<-] (2.5,-5) to [bend right= -10] (1,-3);
                \draw[->-] (2.5,-5) to [bend right= 10] (4,-3);
                \drawddwire{2.5}{-7}{$\catvariableof{2}$}

                \drawhblock{5}{0}{2}{$\bencodingof{\formulaof{2}}$}
                \drawtdwire{4}{-1}{$\catvariableof{2}$}
                \drawddwire{6}{-3}{$\catvariableof{3}$}

            \end{scope}

            \drawtextnode{15}{-4}{$\cdots$}

            \begin{scope}[shift={(15,0)}]
                \drawvariabledot{2.5}{-5}
                \draw[-<-] (2.5,-5) to [bend right= -10] (1,-3);
                \draw[->-] (2.5,-5) to [bend right= 10] (4,-3);
                \drawddwire{2.5}{-7}{$\catvariableof{\catorder\shortminus 2}$}

                \drawhblock{5}{0}{2}{$\bencodingof{\formulaof{\catorder\shortminus 2}}$}
                \drawtdwire{4}{-1}{$\catvariableof{\catorder\shortminus 2}$}
                \drawddwire{6}{-3}{$\catvariableof{\catorder\shortminus 1}$}

            \end{scope}
        }
    \end{example}

    \subsection{Generalization to Runge Kutte methods}

    The generalization to Runge-Kutte methods of higher orders $\seccatorder$ is straight forward:
    Each formula $\formulaof{\catenumerator}$ computing $\catindexof{\catenumerator+1}$ requires further arguments $\catindexof{\catenumerator},\ldots,\catindexof{\catenumerator-\seccatorder+1}$.
    As a consequence, the decomposition of $\hypercoreof{\mathcal{F}}$ involves tensors of larger orders, namely
    \begin{align*}
        \hypercoreofat{\mathcal{F}}{\shortcatvariables,\thirdcatvariable}
        = \contractionof{
            \{\bencodingofat{\formulaof{\catenumerator}}{\catvariableof{\catenumerator+1},\catvariableof{\catenumerator},\ldots,\catvariableof{\catenumerator-\seccatorder+1},\thirdcatvariable}
            \wcols \catenumerator\in[\catorder-1]\}
        }{\shortcatvariables,\thirdcatvariable} \, .
    \end{align*}


    \section{Integration of logical constraints}


    \begin{example}[Constrained parameters in Darcy flows]
        When assumptions are made on the medium variables in Darcy flows, they can be added as further tensors.
    \end{example}

    \subimport{discretized-gaussians/}{discretized-gaussians}


    \section{Noisy observations}

    Let us consider access to observations of the solution function at specific domain elements.
    We capture this by adding tensor cores to each observation.

    When the observations of each values are independent, e.g. discretized Gaussian noise, they are added by elementary networks.
    The corresponding sufficient statistics are $\catindexof{\catenumerator}$ and $(\catindexof{\catenumerator})^2$ to each observed value and the Gaussian noise is a maximum entropy distribution.

    \subsection{Zero-noise limit}

    The zero-noise limit can be investigated based on the characterization of generic maximum entropy distributions \cite{goessmann_tensor_2026-1}.

    \begin{example}[Observations in the toy example]
        In the setting \exaref{exa:EulerToy} let us discretize the solution with $\catorder=2$, i.e. only two values are considered.
        Let us consider observations of both by $\seccatindexof{0}$ and $\seccatindexof{1}$.
        If $\seccatindexof{1}=\seccatindexof{0}+1$, then the zero-noise limit is $\onehotmapofat{\seccatindexof{0},\seccatindexof{1}}{\catvariableof{0},\catvariableof{1}}$.
        The more involved limit is $\seccatindexof{1}\neq\seccatindexof{0}+1$, in which the observations are not agreeing with respect to any solution of the PDE.
        In this case, we can apply the characterization of maximum entropy distributions \cite{goessmann_tensor_2026-1} to characterize the limit.

        If $\seccatindexof{1}=\seccatindexof{0}+2$ the limit is the uniform distribution on the support of the corresponding face, namely
        \begin{align*}
            \frac{1}{2} \cdot \left(
            \onehotmapofat{\seccatindexof{0},\seccatindexof{0}+1}{\catvariableof{0},\catvariableof{1}}
            + \onehotmapofat{\seccatindexof{0}+1,\seccatindexof{0}+2}{\catvariableof{0},\catvariableof{1}}
            \right) \, .
        \end{align*}
    \end{example}

    \bibliographystyle{plainnat}
    \bibliography{../../references.bib}


\end{document}